
Foundation model scaling laws predict performance from model size and dataset size, but treat all data tokens as equivalent. This formulation ignores data quality---a dimension that could substantially affect compute requirements---yet ''quality'' lacks an operational definition suitable for scaling law integration. We propose Q(D), a four-component metric combining deduplication ratio, domain diversity, perplexity score, and token efficiency. On 12 controlled C4 subsets (120GB), we show Q(D) correlates strongly (r = 0.78, p < 0.001) with information density measured via token entropy, n-gram redundancy, and semantic diversity. Deduplication emerges as the strongest predictor (r = 0.72, p = 0.001), explaining 52\% of information density variance. All four components contribute non-redundantly with low measurement variance (coefficient of variation < 10\%). Composite Q(D) explains 61\% of density variance ($R^2$ = 0.61). Our results provide an empirically validated, multi-component quality metric for foundation model pretraining, enabling future research on compute-quality tradeoffs.
